\documentclass[11pt]{article}

\usepackage[margin=1in]{geometry}
\usepackage{amsmath,amssymb,amsthm,mathtools,microtype}
\usepackage[T1]{fontenc}
\usepackage[hidelinks]{hyperref}

\newtheorem{theorem}{Theorem}
\newtheorem{lemma}[theorem]{Lemma}
\newtheorem{corollary}[theorem]{Corollary}

\newcommand{\C}{\mathbb C}
\newcommand{\Eop}{\mathcal E}
\newcommand{\M}{\mathcal M}

\title{A Four-Term Counterexample to the Special Image Conjecture
  in Three Pairs}
\author{Roy van Rijn\\{\small Independent researcher}}
\date{28 July 2026}
\hypersetup{
  pdftitle={A Four-Term Counterexample to the Special Image Conjecture in Three Pairs},
  pdfauthor={Roy van Rijn},
  pdfsubject={An explicit three-pair Image-Mathieu counterexample},
  pdfkeywords={Image Conjecture, Special Image Conjecture,
    Mathieu subspace, SU(2), xz-Conjecture, bihomogenization, contraction}
}

\begin{document}
\maketitle

\begin{abstract}
We give an explicit counterexample to the Special Image Conjecture in
three contraction pairs.  In
\(\C[\tau,w,v,t,z,y]\), let
\[
 f=\tau(t-y)(wz+vt),\qquad g=y.
\]
For the standard contraction map \(\Eop\), we prove for every \(m\geq1\)
that
\[
 \Eop(f^m)=0,\qquad
 [t]\Eop(gf^m)=(-1)^{m-1}(m+1)!\,m!.
\]
Thus every positive power of \(f\) lies in the image of the operators
\(\partial_t-\tau,\partial_z-w,\partial_y-v\), whereas
\(gf^m\notin\M_3\) for every \(m\geq1\).  The polynomial \(f\) has
four terms and the multiplier is linear.  Its two-pair dehomogenization is
exactly Christopher D. Long's counterexample to the Mathieu conjecture for
\(SU(2)\); the present construction is its one-pair bihomogeneous lift.
Consequently the minimum number of pairs in which the Special Image
Conjecture can fail is either two or three.
\end{abstract}

\section{The counterexample}

For \(r\geq1\), write
\[
 R_r=\C[\zeta_1,\ldots,\zeta_r,z_1,\ldots,z_r],
 \qquad
 \M_r=\sum_{i=1}^r(\partial_{z_i}-\zeta_i)R_r.
\]
The Special Image Conjecture asserts that \(\M_r\) is a Mathieu
subspace: if \(h^m\in\M_r\) for every \(m\geq1\), then
\(qh^m\in\M_r\) for every \(q\in R_r\) and all sufficiently large
\(m\); see Zhao~\cite{Zhao2010} and van den Essen, Wright, and
Zhao~\cite{EWZ2011}.

Define the polynomial-valued contraction map
\[
 \Eop_r:R_r\longrightarrow\C[z_1,\ldots,z_r],
 \qquad
 \Eop_r(\zeta^\alpha z^\beta)=\partial_z^\alpha z^\beta.
\]

\begin{lemma}\label{lem:kernel}
For every \(r\geq1\), one has \(\M_r=\ker\Eop_r\).
\end{lemma}

\begin{proof}
For \(h\in R_r\),
\[
 \Eop_r(\partial_{z_i}h)
 =\partial_{z_i}\Eop_r(h)
 =\Eop_r(\zeta_i h).
\]
Thus \(\Eop_r((\partial_{z_i}-\zeta_i)h)=0\), so
\(\M_r\subseteq\ker\Eop_r\).  Conversely,
\[
 \zeta_i h\equiv\partial_{z_i}h\pmod{\M_r}.
\]
Repeatedly applying this congruence gives, for every monomial,
\[
 \zeta^\alpha z^\beta
 \equiv\partial_z^\alpha z^\beta
 =\Eop_r(\zeta^\alpha z^\beta)
 \pmod{\M_r}.
\]
By linearity, \(p\equiv\Eop_r(p)\pmod{\M_r}\).  Since \(\Eop_r\) is
the identity on \(\C[z_1,\ldots,z_r]\), \(\Eop_r(p)=0\) implies
\(p\in\M_r\).
\end{proof}

Use the pairs
\[
 (\tau,t),\qquad(w,z),\qquad(v,y)
\]
and put
\[
 \M_3=(\partial_t-\tau)R_3+
       (\partial_z-w)R_3+
       (\partial_y-v)R_3.
\]

\begin{theorem}\label{thm:main}
Let
\[
 \boxed{\quad
 f=\tau(t-y)(wz+vt),\qquad g=y.
 \quad}
\]
Then, for every \(m\geq1\),
\[
 \Eop_3(f^m)=0,\qquad
 [t]\Eop_3(gf^m)=(-1)^{m-1}(m+1)!\,m!\neq0.
\]
In particular, \(\Eop_3(gf^m)\neq0\).  Consequently
\(f^m\in\M_3\) and \(gf^m\notin\M_3\) for every \(m\geq1\).
\end{theorem}

In expanded form,
\[
\boxed{\;
 f=\tau vt^2-\tau vty+\tau wtz-\tau wyz.
\;}
\]
This displays the advertised four terms directly.

\section{Proof}

First omit the pair \((\tau,t)\).  On
\(\C[W,Z,V,Y]\), define
\[
 \mathcal F(W^aZ^bV^cY^d)
 =\delta_{a,b}\delta_{c,d}\,a!\,c!
\]
and set
\[
 P=(1-Y)(WZ+V),\qquad Q=Y.
\]
Expanding the second factor gives
\[
 P^m=\sum_{k=0}^m\binom mk
 W^{m-k}Z^{m-k}V^k(1-Y)^m.
\]
For a nonzero contraction, the \(V,Y\) pair selects
\([Y^k](1-Y)^m=(-1)^k\binom mk\).  Hence
\begin{equation}\label{eq:pure}
 \mathcal F(P^m)
 =\sum_{k=0}^m(-1)^k\binom mk^2(m-k)!\,k!
 =m!\sum_{k=0}^m(-1)^k\binom mk=0.
\end{equation}
After multiplication by \(Q=Y\), the \(V,Y\) pair instead selects
\([Y^k]Y(1-Y)^m=(-1)^{k-1}\binom m{k-1}\) for \(1\leq k\leq m\).
Therefore
\begin{equation}\label{eq:mixed}
 \begin{aligned}
 \mathcal F(YP^m)
 &=m!\sum_{k=1}^m(-1)^{k-1}\binom m{k-1}\\
 &=(-1)^{m-1}m!.
 \end{aligned}
\end{equation}

Homogenize in the coordinate variables \(Z,Y\), treating \(W,V\) as
coefficient, or dual, variables:
\[
 \widetilde P(t,w,z,v,y)
 =t^2P\left(w,\frac zt,v,\frac yt\right)
 =(t-y)(wz+vt).
\]
Each factor of \(P\) has total \(W,V\)-degree one.  Every monomial of
\(P^m\) that survives \(\mathcal F\) therefore has total
\((Z,Y)\)-degree \(m\), matching its total \(W,V\)-degree.  Its
coordinate homogenization contains
\(t^{2m-m}=t^m\), exactly matching the power \(\tau^m\).
The new contraction contributes \(m!\), and
\eqref{eq:pure} gives
\[
 \Eop_3(f^m)=m!\mathcal F(P^m)=0.
\]

For \(YP^m\), a surviving term of \(P^m\) has total
\((Z,Y)\)-degree \(m-1\).  Its homogenization contains
\(t^{2m-(m-1)}=t^{m+1}\).  Contracting against \(\tau^m\)
leaves one power of \(t\) and contributes \((m+1)!\).  By
\eqref{eq:mixed},
\[
 [t]\Eop_3(yf^m)
 =(m+1)!\mathcal F(YP^m)
 =(-1)^{m-1}(m+1)!\,m!\neq0.
\]
Lemma~\ref{lem:kernel} completes the proof of
Theorem~\ref{thm:main}.

\section{Relation with the
  \texorpdfstring{\(xz\)- and \(SU(2)\)}{xz- and SU(2)} counterexamples}

Write the coordinate functions on \(SU(2)\) as
\[
 \begin{pmatrix}a&c\\ b&d\end{pmatrix}.
\]
Under the substitution
\[
 W=a,\qquad Z=d,\qquad V=b,\qquad Y=-c,
\]
the two-pair seed and multiplier become
\[
 P=(1-Y)(WZ+V)\longmapsto(1+c)(ad+b),\qquad
 Q=Y\longmapsto-c.
\]
This is precisely Long's \(SU(2)\) counterexample~\cite{Long2026}.
Indeed normalized Haar integration satisfies
\[
 \int_{SU(2)}a^r b^s c^t d^u\,dg
 =(-1)^s\delta_{r,u}\delta_{s,t}
   \frac{r!\,s!}{(r+s+1)!}.
\]
Because every monomial of \(P^m\) has total \(W,V\)-degree \(m\),
this formula gives
\[
 \mathcal F(P^m)=(m+1)!\int_{SU(2)}((1+c)(ad+b))^m\,dg
\]
and the analogous identity with multipliers \(Y\) and \(-c\).
Thus equations~\eqref{eq:pure}--\eqref{eq:mixed} recover Long's moments
\[
 \int_{SU(2)}((1+c)(ad+b))^m\,dg=0,\qquad
 \int_{SU(2)}(-c)((1+c)(ad+b))^m\,dg
 =\frac{(-1)^{m-1}}{m+1}.
\]

Under the Müger--Tuset integration substitution used by Long,
\[
 (a,b,c,d)
 =\bigl((1-x)z_2,xz_1,-z_1^{-1},z_2^{-1}\bigr),
\]
the same seed becomes
\[
 (1-z_1^{-1})\bigl((1-x)+xz_1\bigr),
\]
which is Long's three-term counterexample to \(xz(1,1)\).
Accordingly, the seed and its \(xz\)- and \(SU(2)\)-interpretations are
due to Long.  The contribution here is the degree-two one-pair
bihomogenization that produces the four-term counterexample to the Special
Image Conjecture in three pairs.

\begin{corollary}
If \(r_{\mathrm{SIC}}\) denotes the minimum number of contraction pairs
in which the Special Image Conjecture fails, then
\[
 2\leq r_{\mathrm{SIC}}\leq3.
\]
\end{corollary}

\begin{proof}
The upper bound is Theorem~\ref{thm:main}.  The one-pair case follows
from~\cite[Theorem 2.8]{EWZ2011}: take the \(\mathbb Q\)-algebra
\(A=\C[\tau]\) and \(a=\tau\).  Then \(a\) is a non-zero-divisor and
the principal ideal \(Aa=(\tau)\) is radical, precisely the hypotheses
of that theorem.
\end{proof}

\section{Size, scope, and verification}

The coefficients of \(f\) lie in \(\{1,-1\}\).  With respect to
\((\tau,w,v)\)-degree and \((t,z,y)\)-degree, \(f\) has bidegree
\((2,2)\) and ordinary degree four; \(g=y\) has bidegree \((0,1)\).
The polynomial \(f\) is not of the separated form
\(\lambda(\zeta)P(z)\); consequently the example does not give a
three-variable Generalized Vanishing Conjecture counterexample.

Thus the known one-pair result and Theorem~\ref{thm:main} give
\[
 2\leq r_{\mathrm{SIC}}\leq3.
\]
We make no claim here concerning the two-pair case or absolute
literature priority.

The all-order proof above is independent of computation.  An exact
sparse-arithmetic checker expands \(f\), verifies the contractions
through \(m=10\), and checks the two binomial identities through
\(m=99\).  The checker is included in the accompanying Zenodo record
and in the
\href{https://github.com/royvanrijn/jacobian-research}
{companion repository}.\footnote{Repository path:
\path{scripts/verify_three_pair_image_mathieu_counterexample.py}.}
As a low-order checksum,
\[
 \Eop_3(f)=0,\qquad \Eop_3(yf)=2t-y.
\]

\paragraph{Acknowledgment.}
Computational and editorial assistance from OpenAI Codex was used in
searching for, checking, and presenting the example.  The author takes
responsibility for the mathematical statements and proof.

\begin{thebibliography}{9}
\small

\bibitem{Zhao2010}
W.~Zhao,
\emph{Images of commuting differential operators of order one with
constant leading coefficients},
J. Algebra \textbf{324} (2010), 231--247.
\href{https://doi.org/10.1016/j.jalgebra.2010.04.022}
{doi:10.1016/j.jalgebra.2010.04.022}.

\bibitem{EWZ2011}
A.~van den Essen, D.~Wright, and W.~Zhao,
\emph{On the Image Conjecture},
J. Algebra \textbf{340} (2011), 211--224.
\href{https://doi.org/10.1016/j.jalgebra.2011.04.036}
{doi:10.1016/j.jalgebra.2011.04.036}.

\bibitem{Long2026}
C.~D. Long,
\emph{Counterexamples to the \(xz\)-Conjecture and the Mathieu Conjecture
for \(SU(2)\)},
arXiv:2607.19012v1 (2026).
\href{https://arxiv.org/abs/2607.19012}
{arXiv:2607.19012}.

\end{thebibliography}

\end{document}
